% weave_final.tex
% Target: ISSTA 2027 or MSR 2027 workshop track
% All empirical values sourced from RESULTS.md (Phases 5-9b).
% Build: pdflatex weave_final && bibtex weave_final && pdflatex weave_final && pdflatex weave_final

\documentclass[sigconf,review,anonymous]{acmart}

\usepackage{booktabs}
\usepackage{listings}
\usepackage{xcolor}
\usepackage{microtype}
\usepackage{amsmath}
\usepackage{amssymb}
\usepackage{thmtools}
\usepackage{mdframed}

% -------------------------------------------------------
% Theorem-style environment for the formal Definition
% -------------------------------------------------------
\declaretheoremstyle[
  headfont=\bfseries,
  bodyfont=\normalfont,
  mdframed={linewidth=0.5pt, innerleftmargin=6pt, innerrightmargin=6pt,
            innertopmargin=4pt, innerbottommargin=4pt}
]{defstyle}
\declaretheorem[style=defstyle, name=Definition]{definition}

% -------------------------------------------------------
% Go code listing style
% -------------------------------------------------------
\lstdefinelanguage{Go}{
  keywords={func, go, select, case, chan, var, make, close, range, for, return, defer},
  keywordstyle=\bfseries\color{blue!70!black},
  comment=[l]{//},
  commentstyle=\itshape\color{gray},
  stringstyle=\color{green!50!black},
  basicstyle=\ttfamily\small,
  breaklines=true,
}

% -------------------------------------------------------
% ACM metadata — fill before camera-ready
% -------------------------------------------------------
\acmConference[ISSTA/MSR'27]{Workshop on Software Testing and Runtime Analysis}{2027}{TBD}
\acmYear{2027}
\copyrightyear{2027}
\acmISBN{}
\acmDOI{}

% -------------------------------------------------------
\begin{document}

\title{Weave: Distributional Signatures of Concurrent Execution\\
       and Select-Block Goroutine Leak Detection}

\author{Anonymous Author(s)}
\affiliation{\institution{Anonymous Institution}}
\email{anonymous@example.com}

\begin{abstract}
Current execution-trace models treat concurrent programs as if they execute
deterministically. They do not: concurrent programs produce different goroutine
interleavings on every run. We present Weave, a study that exploits this
nondeterminism as a training signal by aggregating multiple runs into
\emph{empirical next-state distributions} and asking frontier language models to
predict distributions rather than point labels.

We make two contributions. \textbf{(1)} On Weave-Bench, a corpus of 26 annotated
concurrent Go programs spanning 13 distinct goroutine leak mechanisms, we show
that distribution framing with explicit reasoning budget reduces Expected
Calibration Error by 21.8\% versus a point-prediction baseline (ECE 0.169 vs.\
0.216), and that model entropy significantly tracks program nondeterminism level
(Spearman $\rho{=}0.412$, $p{=}0.007$). \textbf{(2)} We identify and formally
characterise the \emph{select-block leak}: a class of goroutine leak where a
goroutine enters a permanently blocked \texttt{select} before any \texttt{GoUnblock}
events occur, producing $\mathrm{P(GoUnblock)=0}$ across all trace depths. We
provide a causal proof grounded in Go scheduler semantics, confirm zero false
positives across 26 programs, and verify the signature generalises from 2-case to
4-case \texttt{select} statements. Weave-Bench, all code, and the 377-example
dataset are available at \texttt{https://github.com/kaviru2/Weave}.
\end{abstract}

\keywords{concurrent programs, goroutine leaks, execution traces,
  distributional prediction, calibration, Go}

\maketitle

% =======================================================
\section{Introduction}
% =======================================================

Language models trained or evaluated on execution traces have demonstrated
substantially improved code reasoning in sequential settings~\cite{meta2025cwm}.
The core idea---provide the model with the program state after every executed
line---gives the model a causal model of execution, not just syntax. The Code
World Model (CWM) paper showed this approach dramatically improves bug detection
and next-state prediction for Python programs.

But sequential Python is the easy case. Real systems software is concurrent.
Goroutines communicate over channels and contend on mutexes. The same program can
produce thousands of distinct valid execution orderings. Two fundamental problems
arise. First, point-prediction models produce a single next-event label; this
ignores the distributional nature of the ground truth. Second, existing
benchmarks~\cite{concur2026} confirm that frontier models achieve near-zero
accuracy on concurrent bug detection zero-shot.

We study both problems in the context of Go. Go's \texttt{runtime/trace} package
exposes goroutine scheduler events (\texttt{GoStart}, \texttt{GoBlock},
\texttt{GoUnblock}, \texttt{GoCreate}, \texttt{GoEnd}, \texttt{GoSched}) with
nanosecond timestamps. Running the same program multiple times produces different
event orderings---empirical evidence of the ground-truth distribution over next
states.

\textbf{Contribution 1 (distributional calibration).} We show that asking a
frontier model to predict a probability distribution over the six scheduler event
types, rather than a single event, reduces Expected Calibration Error by 21.8\%
versus the point-prediction baseline, provided the model is given an explicit
reasoning budget. Without reasoning tokens the model produces degenerate near-point
predictions regardless of the distributional prompt. Model entropy with reasoning
enabled significantly correlates with program nondeterminism level
($\rho{=}0.412$, $p{=}0.007$, $n{=}42$), suggesting the model has genuine---if
underexploited---uncertainty about concurrent execution.

\textbf{Contribution 2 (select-block leak signature).} We identify a
formally characterisable class of goroutine leak---the \emph{select-block
leak}---where $\mathrm{P(GoUnblock)=0}$ holds causally, not empirically.
We provide a definition grounded in Go scheduler semantics, confirm zero false
positives across 26 programs (including 11 success, 1 race, and 13 distinct
goroutine-leak mechanisms), and show the signature holds for \texttt{select}
statements with 2 and 4 cases. The key result is that the boundary of the claim is
now precisely known: the signature holds only when the leaked goroutine enters the
\texttt{select} before any \texttt{GoUnblock} events appear in the trace, not for
all goroutine leaks.

% =======================================================
\section{Background}
% =======================================================

\subsection{Code World Models}

Pan et al.~\cite{meta2025cwm} trained a model on Python execution traces (one
state per executed line) and showed dramatically improved code reasoning accuracy
compared to models trained only on source code. The key insight is that the
execution trace provides the causal ground truth that syntax alone does not. We
extend this paradigm to concurrent execution, where the ground truth is a
\emph{distribution} over next states rather than a single state.

Concurrent execution introduces two complications. First, the ground-truth next
state is nondeterministic: even fixing the program and trace prefix, the next
scheduler event can vary across runs. Second, local variable state is not directly
accessible from Go's execution tracer; only goroutine lifecycle events are
available (Section~\ref{sec:bench}).

\subsection{Concurrent Bug Detection}

Tu et al.~\cite{tu2019understanding} studied 171 real concurrency bugs in Go
programs from major open-source projects (Docker, Kubernetes, gRPC, etcd,
CockroachDB). They found that 42\% of bugs are blocking (deadlocks, goroutine
leaks) and 58\% are non-blocking (data races). Their taxonomy directly motivated
our program selection: five programs in Weave-Bench reproduce bug patterns from
their study with direct provenance to specific GitHub issues.

Prior work on static goroutine leak detection uses escape analysis and channel
flow~\cite{locksmith2006} or dynamic analysis tools such as
goleak~\cite{goleak2019}. These approaches produce binary verdicts (leaked or
not). Our distributional approach produces a graded signal across trace depths,
enabling earlier detection from partial traces.

\subsection{Execution Trace Analysis}

TRACED~\cite{traced2024} demonstrated execution-aware pre-training for sequential
code understanding. The CONCUR benchmark~\cite{concur2026} confirms that frontier
models achieve near-zero concurrent bug detection zero-shot, establishing the gap
our work addresses. Debugging CWMs~\cite{debugging_cwm2026} documents known
failure modes when applying world-model training to complex sequential programs;
we expect analogous failure modes to apply in the concurrent setting.

% =======================================================
\section{Weave-Bench}
\label{sec:bench}
% =======================================================

\subsection{Program Corpus}

Weave-Bench contains 26 Go programs spanning four concurrency outcomes
(Table~\ref{tab:programs}). Programs 01--15 were constructed in the initial
corpus phase; programs 16--26 were added specifically to stress-test the
goroutine-leak distribution signature across distinct leak mechanisms.

\begin{table}[h]
\centering
\caption{Weave-Bench corpus summary. ``ASPLOS'' programs reproduce bug patterns
from Tu et al.~\cite{tu2019understanding} with specific GitHub issue provenance.}
\label{tab:programs}
\begin{tabular}{lrrl}
\toprule
\textbf{Outcome} & \textbf{Programs} & \textbf{Groups} & \textbf{Notes} \\
\midrule
success  & 11 & 33 & 11 concurrent patterns \\
leak     & 13 & 39 & 13 distinct mechanisms \\
race     &  1 &  3 & concurrent map write \\
deadlock &  1 &  0 & always times out \\
\midrule
\textbf{Total} & \textbf{26} & \textbf{75} & \\
\bottomrule
\end{tabular}
\end{table}

Five programs reproduce real bug patterns from Tu et al.~\cite{tu2019understanding}
with direct ASPLOS provenance:
\begin{itemize}
  \item \texttt{04\_deadlock}: \texttt{Wait()} inside goroutine-creation loop
        (Docker \#25384)
  \item \texttt{05\_race\_condition}: concurrent map writes without sync
        (ASPLOS non-blocking/shared-memory class)
  \item \texttt{06\_channel\_select}: unbuffered channel + select-timeout leaks
        goroutine (Kubernetes \texttt{finishReq})
  \item \texttt{07\_worker\_pool}: lock acquired after queue pop (Kubernetes quota
        controller)
  \item \texttt{12\_once}: \texttt{sync.Once} prevents channel double-close panic
        (Docker \#24007)
\end{itemize}

Programs 16--25 cover ten additional goroutine-leak mechanisms (HTTP handler,
ticker, worker pool without close, context ignore, event listener, done channel,
mutex wait, pipeline drain failure, select with no default, and per-request
goroutine spawning). Program 26 tests the select-block signature with a
four-case \texttt{select}.

Each program carries a structured \texttt{WEAVE\_META} comment specifying
\texttt{outcome}, \texttt{concurrency\_pattern}, \texttt{goroutine\_count}, and
\texttt{expected\_nondeterminism} (none, low, medium, high).

\subsection{Trace Collection}

We use \texttt{golang.org/x/exp/trace}~\cite{golang_exp_trace} to collect
goroutine scheduler events. Each program is compiled with the race detector
enabled (\texttt{-race}) and run under a 5-second context deadline (500\,ms for
deadlock programs). The trace records six event types: \texttt{GoStart},
\texttt{GoBlock}, \texttt{GoUnblock}, \texttt{GoCreate}, \texttt{GoEnd},
\texttt{GoSched}.

\textbf{Limitation.} The Go runtime trace does not expose local variable state
or heap addresses; only goroutine lifecycle events and function names from stack
frames are available. This is a fundamental constraint of the tracer.

\subsection{Dataset Construction}

Each program is executed five times independently. For each run, we extract three
evaluation examples at trace split points of 25\%, 50\%, and 75\% of total events.
Each example contains the full program source, the partial trace up to the split
point, and the ground-truth next event. Deadlocked programs (outcome \texttt{deadlock})
time out and produce no \texttt{next\_event}; they contribute one example per run
marking \texttt{TimedOut=true}, excluded from accuracy and distribution scoring.

The resulting dataset contains \textbf{377 per-run examples}. For distribution
learning, we aggregate by \texttt{(program\_id, split\_percent)} to form 75
aggregated groups of 5 runs each. We apply a symmetric Dirichlet prior
$\alpha{=}0.5$ (Jeffreys prior~\cite{jeffreys1946}) to obtain posterior
concentration parameters; the empirical distribution is the MLE.

\textbf{Approximation.} Grouping by split percentage treats runs with the same
fractional split point as having comparable trace prefixes. In practice, traces
across runs are structurally similar (the same code paths execute) but not
byte-identical. This approximation is acknowledged and noted as a limitation in
Section~\ref{sec:limitations}.

% =======================================================
\section{Distribution Learning from Empirical Traces}
% =======================================================

\subsection{Empirical Distributions Encode Nondeterminism}

Aggregating five runs per \texttt{(program\_id, split\_percent)} group produces
empirical next-event distributions. Table~\ref{tab:entropy} shows mean Shannon
entropy~\cite{shannon1948} by annotated nondeterminism level, computed over the
42 groups from the original 15-program corpus.

\begin{table}[h]
\centering
\caption{Mean empirical next-event entropy (bits) by annotated nondeterminism
level. Higher nondeterminism correctly produces higher entropy, validating that
empirical distributions carry a real signal.}
\label{tab:entropy}
\begin{tabular}{lrr}
\toprule
\textbf{Nondeterminism} & \textbf{Mean Entropy (bits)} & \textbf{Groups ($n$)} \\
\midrule
high   & 1.396 & 6  \\
medium & 1.210 & 15 \\
low    & 0.903 & 18 \\
none   & 0.971 &  3 \\
\bottomrule
\end{tabular}
\end{table}

Entropy stratifies correctly: high $>$ medium $>$ low. The inversion between
\texttt{none} (0.971) and \texttt{low} (0.903) is within noise at $n{=}3$
\texttt{none}-annotated groups.

\subsection{Deadlock Invisibility}

Program \texttt{04\_deadlock} times out on every run (\texttt{TimedOut=true})
and is excluded from all distribution estimates. The hypothesised
$\mathrm{P(GoBlock)} \to 1$ collapse~\cite{meta2025cwm} cannot be observed for
deadlocks because the tracer never records a \texttt{next\_event}. This is a
fundamental instrumentation limitation: \texttt{trace.Stop()} never executes in a
deadlocked program. We note this limitation and do not claim distribution-based
deadlock detection.

\subsection{Dirichlet Posterior}

Given observed counts $c_k$ for each event type $k \in \{\texttt{GoStart},
\texttt{GoBlock}, \texttt{GoUnblock}, \texttt{GoCreate}, \texttt{GoEnd},
\texttt{GoSched}\}$ across five runs, the Dirichlet posterior under Jeffreys prior
$\alpha{=}0.5$ has concentration parameters:
\[
  \alpha_k^{\mathrm{post}} = 0.5 + c_k
\]
The posterior mean gives a smoothed distribution that assigns non-zero probability
to unobserved events. For the select-block leak signature, the empirical MLE
suffices: $\hat{p}(\texttt{GoUnblock}) = 0$ when $c_{\texttt{GoUnblock}} = 0$
across all five runs.

% =======================================================
\section{Zero-Shot Evaluation}
% =======================================================

\subsection{Point-Prediction Baseline}

We evaluate Gemini 2.5 Flash zero-shot on 207 scored examples from the
15-program corpus (5 deadlock examples excluded). The prompt provides the full
program source, partial execution trace, and current goroutine states; the model
predicts the next event type and goroutine ID.

\begin{table}[h]
\centering
\caption{Zero-shot point-prediction baseline (Gemini 2.5 Flash, $n{=}207$).}
\label{tab:baseline}
\begin{tabular}{lr}
\toprule
\textbf{Metric} & \textbf{Result} \\
\midrule
event\_type accuracy       & 56.0\% (116/207) \\
goroutine\_id accuracy     & 49.8\% (103/207) \\
deadlock detection         & 0/5 \\
race condition detection   & 0/12 \\
high-confidence precision  & 58.0\% \\
\bottomrule
\end{tabular}
\end{table}

\textbf{Failure modes.} The two most common confusions are
\texttt{GoStart}$\leftrightarrow$\texttt{GoUnblock} (20 cases) and
\texttt{GoBlock}$\leftrightarrow$\texttt{GoStart} (19 cases). Block-unblock
symmetry is the dominant failure: the model cannot reliably distinguish which
goroutine transitions or in which direction. Zero deadlock or race detections
confirm that point-prediction framing provides no bug-detection signal.

Accuracy degrades by pattern (fanout 73.3\% vs.\ mutex 43.3\%) and by
nondeterminism level (low 63.3\% vs.\ high 44.4\%), confirming that the model
finds concurrent programs harder as nondeterminism increases.

\textbf{Overconfidence.} High-confidence predictions are correct only 58.0\% of
the time---essentially at chance relative to the model's claimed confidence level.

\subsection{Distribution-Framing Evaluation}

We evaluate Gemini 3.5 Flash on the 42 aggregated groups from the 15-program
corpus, prompting the model to output a probability distribution over all six
event types. We test two conditions: thinking disabled (no reasoning tokens) and
thinking budget of 1,024 tokens. Expected Calibration Error (ECE) is computed as
the mean absolute error between predicted and empirical probabilities, averaged
over six event types.

\textbf{Note on model versions.} The point-prediction baseline used Gemini 2.5
Flash; the distribution evaluation used Gemini 3.5 Flash. The ECE comparison is
therefore not perfectly controlled for model version. We report the values as
measured and acknowledge this confound.

\begin{table}[h]
\centering
\caption{Expected Calibration Error (ECE) comparison. Lower is better. The
point-prediction baseline ECE was recomputed on the full 26-program dataset;
the 21.8\% figure uses this updated baseline.}
\label{tab:ece}
\begin{tabular}{lrr}
\toprule
\textbf{Condition} & \textbf{ECE} & \textbf{vs.\ Baseline} \\
\midrule
Point-prediction (Phase 4, recomputed) & 0.216 & --- \\
Distribution, thinking disabled        & 0.183 & $-$15.3\% \\
Distribution, thinking=1024            & \textbf{0.169} & $-$\textbf{21.8\%} \\
\bottomrule
\end{tabular}
\end{table}

Distribution framing alone (thinking disabled) provides a modest improvement.
The thinking budget is the decisive lever: without it the model produces degenerate
near-point predictions regardless of the distributional prompt (mean entropy
$H{\approx}0$ bits). With thinking enabled, model entropy rises to 0.4--1.0\,bits
across all groups.

\subsection{Entropy Tracks Nondeterminism}

Table~\ref{tab:model_entropy} compares model entropy to empirical entropy by
annotated nondeterminism level.

\begin{table}[h]
\centering
\caption{Mean model entropy (bits) vs.\ empirical entropy by nondeterminism level
(Gemini 3.5 Flash, thinking=1024 vs.\ disabled). Even with thinking enabled the
model systematically underestimates empirical entropy.}
\label{tab:model_entropy}
\begin{tabular}{lrrr}
\toprule
\textbf{Level} & \textbf{No thinking} & \textbf{Thinking=1024} & \textbf{Empirical} \\
\midrule
high   & 0.287 & 1.029 & 1.396 \\
medium & 0.686 & 0.928 & 1.210 \\
low    & 0.042 & 0.399 & 0.903 \\
none   & 0.000 & 0.687 & 0.971 \\
\bottomrule
\end{tabular}
\end{table}

Spearman rank correlation between annotated nondeterminism rank (none=0, low=1,
medium=2, high=3) and model entropy (thinking=1024) over $n{=}42$ groups:
$\rho{=}0.412$, $p{=}0.007$. This is a moderate-to-strong positive correlation
and is statistically significant. With thinking disabled, $\rho{\approx}0$ (no
correlation). The result holds: \emph{thinking budget converts distributional
framing from a formatting choice to genuine uncertainty reasoning}.

The ordering high $>$ medium $>$ low holds; \texttt{none} (0.687) $>$
\texttt{low} (0.399) breaks strict monotonicity, attributable to noise at
$n{=}3$ \texttt{none}-annotated groups.

\textbf{KL divergence.} Mean KL(empirical $\|$ predicted) drops from 9.50 nats
(thinking disabled) to 7.00 nats (thinking=1024). The residual gap of 7.00 nats
quantifies the training target for a model optimised to minimise KL from
empirical distributions.

\subsection{Entropy Decreases with Trace Depth}

Table~\ref{tab:entropy_depth} shows model entropy decreases monotonically as the
model observes a larger fraction of the trace. Empirical entropy remains stable
across depths (${\approx}1.1$ bits), confirming that the model's increasing
confidence is from reasoning over trace context, not a property of the programs
themselves at that depth.

\begin{table}[h]
\centering
\caption{Model and empirical entropy (bits) by trace split depth.
Spearman $\rho{=}-0.314$, $p{=}0.043$ (depth vs.\ model entropy, $n{=}42$).}
\label{tab:entropy_depth}
\begin{tabular}{lrr}
\toprule
\textbf{Trace depth} & \textbf{Model entropy} & \textbf{Empirical entropy} \\
\midrule
25\% & 0.941 & 1.099 \\
50\% & 0.709 & 1.113 \\
75\% & 0.445 & 1.051 \\
\bottomrule
\end{tabular}
\end{table}

The model entropy-depth trend is statistically significant ($p{=}0.043$),
showing the model does integrate trace information into its uncertainty
estimates. The stable empirical entropy confirms this is a model behaviour, not a
dataset artefact.

\subsection{Anomaly Scores (Underpowered)}

We define an anomaly score $\mathrm{KL}(\hat{p} \| \mathrm{Uniform})$ for each
group: high values indicate the model is confident (distributional mass is
concentrated), low values indicate high uncertainty (novel or unusual execution
pattern). Mean anomaly scores: success programs 1.342, buggy programs (leak +
race) 1.182. Cohen's $d{=}0.294$ (small effect), $p{=}0.503$---not statistically
significant at $n{=}9$ buggy groups. We report this result honestly: the anomaly
signal exists but is too weak to use in practice at this dataset scale.

% =======================================================
\section{Select-Block Goroutine Leak Signature}
\label{sec:leak}
% =======================================================

\subsection{Motivation}

In the original 15-program corpus, both goroutine-leak programs
(\texttt{06\_channel\_select} and \texttt{14\_goroutine\_leak}) showed
$\mathrm{P(GoUnblock)=0}$ across all trace split depths, while all success
programs showed $\mathrm{P(GoUnblock)} > 0$ at some depth. This suggested a
distribution-based leak detector. Before claiming this as a contribution, we
stress-tested the signature across 11 additional leak mechanisms.

\subsection{Formal Definition}

We first define the signature class formally.

\begin{definition}[Select-block leak]
\label{def:selectblock}
A goroutine $G$ is a \emph{select-block leak} if there exists a time $t$ at
which $G$ enters a \texttt{select} statement where no case condition can be
satisfied by any reachable future execution of the program. For all $t' > t$:
$\mathrm{GoUnblock}(G)$ is structurally impossible.

Consequently, $\mathrm{P(GoUnblock)=0}$ holds in the empirical next-event
distribution for any trace split point that falls after $G$'s first \texttt{GoBlock}
event at time $t$.
\end{definition}

\textbf{Causal justification.} In the Go runtime, \texttt{GoUnblock($G$)} fires
only when another goroutine sends to or closes a channel on which $G$ is blocked,
or unlocks a mutex $G$ is waiting on~\cite{gopkg_trace}. If no such action is
reachable from the remaining program execution, \texttt{GoUnblock} cannot occur.
This is a theorem about Go scheduler semantics, not an empirical observation.

\subsection{Empirical Confirmation}

\begin{table}[h]
\centering
\caption{$\mathrm{P(GoUnblock)}$ by outcome and trace split depth.
Values are means over groups at each depth. Success programs maintain
non-zero $\mathrm{P(GoUnblock)}$ at all depths.}
\label{tab:leak_sig}
\begin{tabular}{llrrr}
\toprule
\textbf{Outcome} & \textbf{Programs} & \textbf{25\%} & \textbf{50\%} & \textbf{75\%} \\
\midrule
success  & 11 & 0.182 & 0.091 & 0.218 \\
leak (all 13) & 13 & 0.169 & 0.077 & 0.246 \\
race    &  1 & 0.000 & 0.000 & 0.000$^*$ \\
\midrule
\multicolumn{5}{l}{\small select-block programs only:} \\
\texttt{06\_channel\_select}      & & 0.000 & 0.000 & 0.000 \\
\texttt{24\_select\_no\_default}  & & 0.000 & 0.000 & 0.000 \\
\texttt{26\_select\_block\_multicase} & & 0.000 & 0.000 & 0.000 \\
\bottomrule
\multicolumn{5}{l}{$^*$See Section~\ref{sec:race_confound}.}
\end{tabular}
\end{table}

Expanding from 2 to 13 goroutine-leak programs reveals that $\mathrm{P(GoUnblock)=0}$
across \emph{all} split depths holds only for the three select-block leak programs
(Table~\ref{tab:leak_sig}). The aggregate mean over all 13 leak programs shows
non-zero $\mathrm{P(GoUnblock)}$, confirming the signature is mechanism-dependent.

The other 10 leak mechanisms---workers waiting on unclosed channels, goroutines
ignoring context cancellation, goroutines awaiting mutex locks, etc.---all produce
$\mathrm{P(GoUnblock)} > 0$ at some split depth because their goroutines perform
legitimate work (receiving items, processing requests) before reaching the
permanently blocked state. GoUnblock events generated during this legitimate
work appear in early split windows.

\textbf{Zero false positives.} No success program shows $\mathrm{P(GoUnblock)=0}$
across all three split depths. The select-block signature has zero false positives
at the 26-program scale of this study.

\subsection{Multi-Case Boundary Test}

Program \texttt{26\_select\_block\_multicase} tests whether the signature holds
for selects with more than two cases. The program contains a four-case
\texttt{select} over channels of types \texttt{int}, \texttt{string},
\texttt{[]byte}, and \texttt{error}, none of which are ever sent to.

\begin{lstlisting}[language=Go, caption={Four-case select-block leak
(program 26). The goroutine blocks in \texttt{GoWaiting} immediately;
no \texttt{GoUnblock} events can occur.}]
go func() {
    select {
    case v := <-intCh:    fmt.Println("int:", v)
    case s := <-strCh:    fmt.Println("string:", s)
    case b := <-byteCh:   fmt.Println("bytes:", b)
    case err := <-errCh:  fmt.Println("error:", err)
    }
}()
\end{lstlisting}

Result: $\mathrm{P(GoUnblock)=0}$ at all three split depths across all five runs.
The select arity does not affect the signature. What matters is the structural
property in Definition~\ref{def:selectblock}: no case is reachable, so
\texttt{GoUnblock} is impossible regardless of how many cases the \texttt{select} has.

\subsection{Race Condition Confound}
\label{sec:race_confound}

Program \texttt{05\_race\_condition} (outcome: \texttt{race}) also shows
$\mathrm{P(GoUnblock)=0}$ across all split depths, but for an unrelated reason.
The bug is concurrent writes to a shared map without synchronisation; the program
uses no channels or mutex primitives in its bug path. The WaitGroup unlock---the
only source of \texttt{GoUnblock} events in this program---fires at the very end
of a short trace, consistently falling outside all three split windows.

This is a \textbf{trace-structure confound}, not a consequence of the race
condition. We state this explicitly: the $\mathrm{P(GoUnblock)=0}$ observation
for program 05 is not evidence of the select-block signature, and race condition
detection via this signature is not claimed.

% =======================================================
\section{Limitations}
\label{sec:limitations}
% =======================================================

\textbf{Mechanism-dependent leak signature.} $\mathrm{P(GoUnblock)=0}$ is not a
general goroutine-leak detector. It holds only for the select-block class
(Definition~\ref{def:selectblock}) where the goroutine blocks on an unreachable
\texttt{select} before any legitimate goroutine lifecycle work begins. Eleven of
13 leak mechanisms in Weave-Bench do not satisfy this condition.

\textbf{Deadlock invisibility.} Programs that deadlock produce no trace output
because \texttt{trace.Stop()} never executes. The hypothesised
$\mathrm{P(GoBlock)} \to 1$ collapse for deadlock programs cannot be observed.
Deadlock detection via empirical distributions is not feasible with the current
instrumentation.

\textbf{No local variable state.} The Go execution tracer exposes goroutine
lifecycle events and function names from stack frames but not local variable
values or heap object addresses. Channel and mutex object identities are not
available. This limits the state representation compared to sequential trace
models~\cite{meta2025cwm}.

\textbf{Split-percent approximation.} Grouping runs by split percentage treats
structurally similar (but not identical) trace prefixes as equivalent. The
25\%/50\%/75\% split points of different runs index different absolute events.
This is an approximation that should be replaced by semantic prefix matching
in future work.

\textbf{Underpowered anomaly detection.} The KL-based anomaly score shows a
small but non-significant separation between success and buggy programs (Cohen's
$d{=}0.294$, $p{=}0.503$, $n_{\mathrm{buggy}}{=}9$). At the current dataset
scale this result is inconclusive.

\textbf{Model version mismatch.} The point-prediction baseline used Gemini 2.5
Flash; the distribution evaluation used Gemini 3.5 Flash. The ECE comparison
across these conditions is not fully controlled. A clean comparison would hold
the model constant across both evaluations.

\textbf{Dataset scale.} Weave-Bench contains 26 programs and 377 examples. This
is a pilot benchmark---sufficient to establish feasibility and discover the
select-block signature, but insufficient for statistically robust claims about
anomaly detection or to serve as a held-out test set for trained models. A
conference submission would require 50+ programs.

% =======================================================
\section{Related Work}
% =======================================================

\textbf{Execution trace models.} Pan et al.~\cite{meta2025cwm} is the direct
antecedent: Code World Models train on Python execution traces (state after every
line) and substantially improve code reasoning. Our work extends this paradigm to
concurrent execution, reformulating the prediction target as a distribution.
TRACED~\cite{traced2024} provides execution-aware pre-training for source code
understanding in sequential settings; the concurrent analogue remains open.

\textbf{Concurrent bug benchmarks.} Tu et al.~\cite{tu2019understanding} provide
the most comprehensive taxonomy of Go concurrency bugs in production code and
directly motivate five programs in Weave-Bench. The CONCUR
benchmark~\cite{concur2026} evaluates frontier models on concurrent program
understanding and confirms near-zero concurrent bug detection zero-shot,
establishing the baseline gap our work addresses.

\textbf{Goroutine leak detection.} Static analysis tools (e.g., goroutine escape
analysis) and dynamic tools such as goleak~\cite{goleak2019} detect goroutine
leaks at program exit. These produce binary verdicts; our approach produces a
graded distributional signal observable from partial traces, enabling earlier
detection and uncertainty quantification.

\textbf{Debugging world models.} Debugging CWMs~\cite{debugging_cwm2026} studies
failure modes of execution-trace models in sequential settings; Section~3 of that
paper identifies classes of programs where trace-based models fail, providing
context for the analogous limitations we document in the concurrent setting.

% =======================================================
\section{Conclusion}
% =======================================================

We presented Weave, a study of distributional next-state prediction for concurrent
Go programs. Our two main contributions are:

\textbf{(1) Distributional calibration.} Asking a frontier language model to
predict a probability distribution over scheduler events, combined with an explicit
reasoning budget, reduces ECE by 21.8\% versus a point-prediction baseline and
produces model entropy that significantly tracks program nondeterminism
($\rho{=}0.412$, $p{=}0.007$). The thinking budget is necessary: without it the
model produces degenerate near-point predictions. This reveals latent uncertainty
knowledge that standard prompting suppresses.

\textbf{(2) Select-block leak signature.} We formally defined the select-block
goroutine leak class (Definition~\ref{def:selectblock}), provided a causal proof
grounded in Go scheduler semantics, and confirmed $\mathrm{P(GoUnblock)=0}$
across all trace depths for three programs with 2-case and 4-case
\texttt{select} statements. The signature has zero false positives across 26
programs. We precisely characterised the boundary: the signature holds only when
the leaked goroutine enters the \texttt{select} before any \texttt{GoUnblock}
events occur.

\textbf{Next steps.} The empirical work establishes two concrete training targets:
(i) minimise KL divergence from empirical next-event distributions (the 7-nat
gap), and (ii) use P(GoUnblock) as a distributional feature for select-block leak
classification. Both require training on larger trace corpora, which we plan to
pursue with WSO2 infrastructure on the Ballerina concurrent language.

\bibliographystyle{ACM-Reference-Format}
\bibliography{weave_final}

\end{document}
